\section{Connection to the project's spatial free-surface model}
\label{sec:project}
This appendix identifies a later water-wave convention; it does not replace
the original paper's formulation or change the project's approved model.

The 1979 paper uses a temporal envelope, includes two horizontal modulation
directions, and introduces $A$ through the velocity potential. A wave-tank
experiment may instead prescribe a temporal profile at an upstream location
and ask how it changes with propagation distance. This motivates a
\emph{spatial} evolution formulation.

Your current project contract identifies a one-coordinate free-surface
envelope $u(\xi,\tau)$ with
\begin{equation}\label{eq:spatial-model}
 u_\xi=-iu_{\tau\tau}-i|u|^2u
 -\varepsilon\left(8|u|^2u_\tau+2u^2u_\tau^*
                  +2iu|D_\tau|(|u|^2)\right).
\end{equation}
Its contract follows the free-surface equation (1.3) in \cite{fedele}.
Here $\xi$ is propagation distance and $\tau$ is retarded time along the
profile. In that source's carrier-scaled coordinates,
\[
 B=\varepsilon u,\qquad
 \tau=\varepsilon(2x-t),\qquad \xi=\varepsilon^2x.
\]
The companion potential-envelope equation has a different coefficient for
the conjugate-gradient term. Equality of leading envelope magnitudes does
not justify replacing one by the other at fourth order.

The operator $|D_\tau|$ uses Fourier multiplier $|\kappa|$, with
$e^{-i\kappa\tau}$ as the forward transform kernel and zero derivative
response for the constant mode. It expresses an eliminated mean-flow
response, just as the half-space solution did in Section~\ref{sec:mean}.
It is not obtained by simply renaming the $x$ coordinate in (2.19).

Two array axes $(\xi,\tau)$ describe an evolution coordinate and one
profile coordinate. They do not represent the two horizontal modulation
directions $(x,y)$ of the original stability problem. Comparing models
requires tracking dependent variable, scaling, direction count,
evolution variable, and nonlocal-operator convention.

For learning this paper, retain its notation through the stability
derivation. Afterwards, use this appendix to recognize why a later
water-wave solver can contain related effects with different-looking
coefficients. A numerical implementation should use its separately verified
equation contract rather than this pedagogical transcription of the scan.
\clearpage
